\section{Constructive and Destructive Regimes}
\label{sec:displacement}

Loss spikes and oscillatory updates motivate diagnostics that explain how
training is progressing as it happens. The gradients already computed by
the optimizer contain two complementary signals: how strongly the loss
changes locally, and whether the next update reinforces or opposes the
preceding one. We first show how these interactions determine accumulated
parameter change. We then connect consecutive alignment to curvature,
obtaining an inexpensive estimator and a joint condition on gradient
reversal and growth that explains when the loss can increase.

Consider the plain SGD update in Eq.~\eqref{eq:sgd_map_stability}, with
constant learning rate $\eta>0$. Summing the updates gives the displacement
from initialization,
\begin{equation}
  \Delta_T:=\theta_T-\theta_0=-\eta\sum_{t=1}^T g_t.
  \label{eq:displacement}
\end{equation}
For training from scratch, $\theta_0$ is the initial model and $\theta_T$
contains the parameter changes accumulated while learning from the data.
Expanding their squared difference separates individual step magnitudes
from interactions across steps:
\begin{equation}
\begin{aligned}
  \|\Delta_T\|^2
  ={}&\eta^2\sum_{t=1}^T\|g_t\|^2\\
     &+2\eta^2\sum_{1\leq u<t\leq T}g_u^\top g_t.
\end{aligned}
\label{eq:magnitude-interaction}
\end{equation}
The squared norms are nonnegative contributions from individual steps.
In full-batch GD, the negative gradient is the direction of steepest
local descent, and a larger norm means a steeper local slope. The
cross-step terms determine how these changes accumulate. Positive inner
products reinforce displacement; negative inner products produce
\emph{destructive interference}, with one update undoing a component
of another. Large individual steps can therefore coexist with limited
net parameter change.

To make the respective roles of magnitude and direction explicit,
factor each inner product as
$g_u^\top g_t=\|g_u\|\|g_t\|\rho_{u,t}$, where
\begin{equation}
  \rho_{u,t}:=\frac{g_u^\top g_t}{\|g_u\|\|g_t\|},
  \qquad 1\leq u<t\leq T.
  \label{eq:rho}
\end{equation}
For nonzero consecutive gradients, write $\rho_t:=\rho_{t-1,t}$.
Substitution into Eq.~\eqref{eq:magnitude-interaction} gives
\begin{equation}
\begin{aligned}
\|\Delta_T\|^2={}&\eta^2\sum_{t=1}^T\|g_t\|^2\\
 &+2\eta^2\sum_{1\leq u<t\leq T}
    \|g_u\|\|g_t\|\rho_{u,t}.
\end{aligned}
\label{eq:magnitude-alignment}
\end{equation}
Constructive and destructive behavior are thus encoded directly in the
SGD update equations. For fixed norms, positive alignment increases net
displacement, whereas negative alignment limits it through cancellation.
We identify a \emph{constructive regime} by persistent $\rho_t>0$ and
a \emph{destructive regime} by persistent $\rho_t<0$.

\subsection{Directional Curvature and Consecutive Alignment}

We now connect these regimes to the curvature encountered along a step.
For mini-batch SGD, use
$g_t=\nabla\mathcal L(\theta_{t-1})+\xi_t$ from
Eq.~\eqref{eq:sgd_noise_decomposition}. Suppose $\mathcal L$ is twice
continuously differentiable along the preceding step and write
$H_{t-1}:=\nabla^2\mathcal L(\theta_{t-2})$.
Expanding the full gradient around $\theta_{t-2}$ gives
\begin{equation}
  g_t=(I-\eta H_{t-1})g_{t-1}+(\xi_t-\xi_{t-1})+r_t,
  \label{eq:gradient-evolution}
\end{equation}
where $r_t$ is the Taylor remainder. Neglecting it and projecting onto
the preceding gradient yields
\begin{equation}
\begin{aligned}
\|g_{t-1}\|\|g_t\|\rho_t
 &\approx\underbrace{\|g_{t-1}\|^2}_{\text{Gradient magnitude}}\\
 &\quad-\underbrace{\eta g_{t-1}^{\top}H_{t-1}g_{t-1}}_{\text{Curvature contribution}}\\
 &\quad+\underbrace{g_{t-1}^{\top}(\xi_t-\xi_{t-1})}_{\text{Noise contribution}}.
\end{aligned}
\label{eq:rho-expanded}
\end{equation}
Define the curvature along the preceding gradient and the normalized
noise contribution as
\begin{equation}
  \kappa_t:=\frac{g_{t-1}^\top H_{t-1}g_{t-1}}{\|g_{t-1}\|^2},
  \label{eq:kappa}
\end{equation}
\begin{equation}
  \varepsilon_t:=
  \frac{g_{t-1}^\top(\xi_t-\xi_{t-1})}{\|g_{t-1}\|^2}.
  \label{eq:epsilon}
\end{equation}
Dividing Eq.~\eqref{eq:rho-expanded} by the two gradient norms gives
the central relation
\begin{equation}
  \boxed{\rho_t\approx\frac{\|g_{t-1}\|}{\|g_t\|}
                   (1-\eta\kappa_t+\varepsilon_t).}
  \label{eq:rho-normalized}
\end{equation}
In full-batch GD, $\varepsilon_t=0$ and the norm ratio is positive.
Within this approximation, $\rho_t\propto1-\eta\kappa_t$, so
$\eta\kappa_t<1$ gives constructive alignment, while $\eta\kappa_t>1$
gives destructive alignment. The boundary $\eta\kappa_t=1$ separates
the two regimes. At this boundary the gradients are orthogonal if
both are nonzero; a vanishing gradient leaves the cosine undefined.

This relation supports two interpretations. From a representation-learning
perspective, a destructive update opposes parameter changes made by its
predecessor, providing a mechanism for repeatedly overwriting learned
structure. Even with iid data from a fixed distribution, successive
updates can repeatedly undo previously accumulated parameter changes.
From an optimization perspective, negative alignment indicates
that a GD step has passed the minimum along its preceding direction in
the local quadratic model. The optimizer then partially reverses that
step. Such reversals can coexist with convergence; the magnitude of the
returning gradient determines when they instead increase the loss.

\subsection{Estimating Curvature without Hessian Computation}
\label{subsec:curvature-diagnostic}

Equation~\eqref{eq:rho-normalized} turns consecutive training gradients
into a directional-curvature estimator:
\begin{equation}
  \eta\widehat\kappa_t
  :=1-\rho_t\frac{\|g_t\|}{\|g_{t-1}\|}
   \approx\eta\kappa_t-\varepsilon_t.
\label{eq:kappa-estimate}
\end{equation}
Its computation requires retaining one additional gradient and taking
norms and an inner product. Thus curvature information can be monitored at every
training step using gradients the optimizer already needs.

The measured projected change has a precise geometric meaning. In smooth
full-batch GD, $\widehat\kappa_t$ is exactly the directional Hessian
curvature averaged along the preceding step; it approximates the
starting-point Rayleigh quotient $\kappa_t$ when curvature varies little
over that step. Appendix~\ref{appendix:loss-curvature} gives the identity
and an error bound. It measures curvature along the realized gradient
direction, while $\lambda_{\max}(H_t)$ measures the sharpest direction.
In SGD, the additional term $-\varepsilon_t$ incorporates sampling.
Larger batches reduce the gradient sampling variance and can
bring the diagnostic closer to its full-batch interpretation. Its
importance still depends on gradient variability relative to the full
gradient, which we examine next.

\subsection{Noise Contribution under Mini-Batch Sampling}
\label{subsec:sampling-noise}

Fresh, conditionally unbiased batches satisfy
$\mathbb E[g_{t-1}^\top\xi_t]=0$ and
$\mathbb E[\nabla\mathcal L(\theta_{t-2})^\top\xi_{t-1}]=0$.
Consequently, the noise contribution to consecutive alignment has
nonpositive expectation:
\begin{equation}
  \mathbb E[g_{t-1}^\top(\xi_t-\xi_{t-1})]
  =-\mathbb E[\|\xi_{t-1}\|^2]\leq0.
  \label{eq:noise-expectation}
\end{equation}
For uniform batches sampled without replacement,
\begin{equation}
  \mathbb E[\|\xi_{t-1}\|^2]
  =\frac{N-b}{b(N-1)}
    \mathbb E[\operatorname{tr}\Sigma(\theta_{t-2})].
  \label{eq:finite-population-noise}
\end{equation}
Smaller batches therefore increase the magnitude of this negative
unnormalized contribution at fixed parameters; full-batch evaluation
eliminates it.

For two independent gradients sampled at the same fixed state $\theta$,
concentration of their norms gives the approximation
\begin{equation}
  \mathbb E[\rho_t\mid\theta\text{ fixed}]\approx
  \frac{\|\nabla\mathcal L(\theta)\|^2}
  {\|\nabla\mathcal L(\theta)\|^2+
   \frac{N-b}{b(N-1)}\operatorname{tr}\Sigma(\theta)}.
  \label{eq:noise-cosine-approx}
\end{equation}
Reducing $b$ weakens this shared directional signal, moving the expected
cosine toward zero. Along an evolving SGD trajectory, sampling also
changes the next parameter state and its curvature. The negative term
in Eq.~\eqref{eq:noise-expectation} alone therefore does not determine
how a negative cosine changes. Section~\ref{sec:sgd-noise-experiments}
tests the attenuation of constructive and destructive alignment on
actual training trajectories.

\subsection{Normalized Displacement and Cumulative Interference}
\label{subsec:normalized-displacement}

How much of the distance traveled during training remains as net
parameter change? We measure this fraction by the normalized displacement
\begin{equation}
  R_T:=
  \frac{\|\Delta_T\|}
       {\eta\sum_{t=1}^T\|g_t\|}
  =\frac{\|\sum_{t=1}^T g_t\|}
        {\sum_{t=1}^T\|g_t\|},
  \qquad 0\leq R_T\leq1,
  \label{eq:normalized-displacement}
\end{equation}
defined for nonzero total path length. For constant-step plain SGD, the
denominator is the traveled distance and the numerator is displacement
from initialization. Equation~\eqref{eq:magnitude-alignment} shows that
$R_T^2$ incorporates the interactions between every pair of steps.
Thus $R_T$ summarizes the cumulative effect of their magnitudes and
alignment, while $\rho_t$ identifies the local regime.

Aligned steps give $R_T=1$. Equal opposing pairs give zero, and
equal-length orthogonal steps give $R_T=1/\sqrt T$. Even without direct
cancellation, orthogonal steps retain only 10\% of their traveled
distance after 100 updates, or 1\% after 10,000. The example shows how
weak directional reinforcement limits net movement relative to path
length. Low $R_T$ therefore reveals repeated reversal or weak coherence;
its effect on learning is assessed jointly with the loss.

\input{sections/04-controlled-figures}

\subsection{Conditions for Loss Decrease and Spike Formation}
\label{subsec:loss-connection}

When does a destructive regime produce a loss increase? Define the
full-objective change over the preceding step as
$\Delta\mathcal L_{t-1}:=\mathcal L(\theta_{t-1})-
\mathcal L(\theta_{t-2})$.
The second-order approximation derived in Appendix~\ref{appendix:loss-curvature}
separates three contributions:
\begin{equation}
\begin{aligned}
\Delta\mathcal L_{t-1}
&\approx\underbrace{-\frac{\eta}{2}\|g_{t-1}\|^2}_{\text{Gradient magnitude}}\\
&\quad\underbrace{-\frac{\eta}{2}\rho_t\|g_{t-1}\|\|g_t\|}_{\text{Directional interference}}\\
&\quad+\underbrace{\frac{\eta}{2}g_{t-1}^\top(\xi_{t-1}+\xi_t)}_{\text{Noise interaction}}.
\end{aligned}
\label{eq:loss-curvature}
\end{equation}
The first term favors descent. Constructive alignment adds a further
decrease, whereas destructive alignment offsets part of it. The size
of this opposing contribution depends jointly on reversal and gradient
growth. In full-batch GD with nonzero gradients, the approximation
predicts $\Delta\mathcal L_{t-1}>0$ precisely when
\begin{equation}
  \boxed{\rho_t<-\frac{\|g_{t-1}\|}{\|g_t\|}.}
  \label{eq:loss-rise-alignment}
\end{equation}
Two conditions are necessary: $\rho_t<0$ and
$\|g_t\|>\|g_{t-1}\|$. The new gradient must reverse direction strongly
enough that its opposing component exceeds the preceding magnitude.
For example, doubling the norm requires $\rho_t<-0.5$, an angle above
$120^\circ$. A fivefold increase requires only $\rho_t<-0.2$, or about
$101.5^\circ$. Greater norm growth therefore makes even a milder
reversal sufficient for an upward loss step.

The equivalent curvature form is
\begin{equation}
  \Delta\mathcal L_{t-1}\approx
  \frac{\eta}{2}\|g_{t-1}\|^2(\eta\kappa_t-2).
  \label{eq:loss-rise-curvature}
\end{equation}
It summarizes the same joint effect: $\eta\kappa_t=1$ marks the
constructive--destructive boundary, while $\eta\kappa_t=2$ marks
the loss-increase boundary. Between them, destructive updates still
decrease the objective. The expressions are exact for quadratic
objectives; for nonlinear losses their accuracy depends on variation
along the step, with sampling adding the explicit noise contribution.
A sequence of upward steps followed by recovery forms the spike
events studied below. The diagnostic thus connects local curvature
to the reversal and norm growth underlying their formation.
